\documentclass[fontsize=16pt]{mybib}


\graphicspath{{../../../assets/images/}}
\setheaderlogo{mnr-green}
\setdefault{SCHL}
\setversions{SCHL}
\begin{document}

    \section{Bibeltext}
    \begin{block}[\bibleverse{Joh}(12:12-19)]
        \verstab{12}{SCHL}{
            Am folgenden Tag, \konj[b, p]{als} viele \person[b, p]{Leute}, die zum Fest \verbN[b, p]{erschienen waren}, \verbN[b, p]{hörten}, \konj[b, p]{dass} \person[b, p]{Jesus} nach \ort[b, p]{Jerusalem} \verbN[b, p]{komme},
        }
        \begin{block}
            \begin{itemize}
                \item \annot{Am folgenden Tag}: Der Tag, nachdem Maria Jesus die Füsse gesalbt hat.
                \item \annot{als viele Leute die zum Fest erschienen waren}: Das Passah fest. Wie in \bibleverse{Joh}(12:1) zu lesen, war es 6 Tage vor dem Passah Fest.
                \item \annot{, hörten, dass:} Viele der Menschen haben vernommen, dass Jesus nach Jerusalem kommt. Viele kennen seine wundertaten, die Jesus gemacht hat, und wollten diesen Heiler sehen.
            \end{itemize}
        \end{block}
        \verstab{13}{SCHL}{
            da \verbN[b, p]{nahmen} sie Palmzweige und \verbN[b, p]{gingen} hinaus, ihm \verbN[b, p]{entgegen}, und \verbN[b, p]{riefen}: Hosianna! Gepreisen \verbN[b, p]{sei} der, welcher \verbN[b, p]{kommt} im Namen des \person[b, p]{Herrn}, der \person[b, p]{König} von \ort[b, p]{Israel}!
        }
        \begin{block}
            \annot{nahmen sie Palmenzweige:} In der Region gibt es viele Palmenzeweige, mit den Palmzweigen bezeugt die Menge ihre Freude und markieren so den Einzug des Königs. 
        \end{block}
        \verstab{14}{SCHL}{
            \person[b, p]{Jesus} \verbN[b, p]{hatte} \konj[b, p]{aber} einen jungen Esel \verbN[b, p]{gefunden} und \verbN[b, p]{setzte} sich darauf, wie \verbN[b, p]{geschrieben steht}:
        }
        \verstab{15}{SCHL}{
            \verbI[b, p]{Fürchte} dich nicht, Tochter Zion! \verbI[b, p]{Siehe}, dein \person[b, p]{König} \verbN[b, p]{kommt}, \verbP[b, p]{sitzend} auf dem Füllen einer Eselin.
        }
        \verstab{16}{SCHL}{
            Dies \verbN[b, p]{verstanden} \konj[b, p]{aber} seine \person[b, p]{Jünger} anfangs nicht, \konj[b, p]{doch} als \person[b, p]{Jesus} \verbN[b, p]{verherrlicht war}, \konj[b, p]{da} erinnerten sie sich, \konj[b, p]{dass} dies von ihm \verbN[b, p]{geschrieben} stand und \konj[b, p]{dass} sie ihm dies \verbN[b, p]{getan hatten}.
        }
        \verstab{17}{SCHL}{
            Die \person[b, p]{Menge} nun, die bei ihm \verbN[b, p]{war}, \konj[b, p]{als} er \person[b, p]{Lazarus} aus dem Grab \verbN[b, p]{gerufen} und ihn aus den Toten \verbN[b, p]{auferweckt hatte}, \verbN[b, p]{legte} Zeugnis \verbN[b, p]{ab}.
        }
        \verstab{18}{SCHL}{
            \konj[b, p]{Darum} \verbN[b, p]{ging} ihm auch die \person[b, p]{Volksmenge} \verbN[b, p]{entgegen}, \konj[b, p]{weil} sie \verbN[b, p]{gehört hatten}, \konj[b, p]{dass} er diese Zeichen \verbN[b, p]{getan hatte}.
        }
        \verstab{19}{SCHL}{
            \konj[b, p]{Da} \verbN[b, p]{sprachen} die \person[b, p]{Pharisäer} zueinander: Ihr \verbN[b, p]{seht}, \konj[b, p]{dass} ihr nichts \verbN[b, p]{ausrichtet}. \verbN[b, p]{Siehe}, alle Welt \verbN[b, p]{läuft} ihm \verbN[b, p]{nach}!
        }
    \end{block}


    

\end{document}